%----------------Cap_05----------------%
\chapter{Conclusão}
\label{cap:conclusao}

Este trabalho apresentou o projeto, a implementação e a validação experimental de uma arquitetura de controle preditivo neural embarcado (\textit{TinyML}), concebida para aproximar a política de um Controlador Preditivo Não-Linear (NMPC) ótimo e executá-la em tempo real em um microcontrolador de baixo custo para o rastreamento angular de um aeropêndulo instrumentado.

A motivação central da pesquisa fundamentou-se na superação do gargalo computacional associado à otimização numérica não-linear online: a resolução iterativa do problema de controle ótimo via resolvedores tradicionais (como o IPOPT no ambiente CasADi) exigiu tempos de convergência entre 90 e 460\,ms por ciclo de amostragem em computador de mesa. Diante de uma dinâmica mecatrônica rápida que impõe um período de controle de $T_s = 10$\,ms ($100$\,Hz), a transferência da inteligência preditiva para uma rede neural compacta em linguagem C viabilizou a execução em hardware embarcado com determinismo temporal estrito e elevada concordância dinâmica frente à lei ótima nominal.

% ============================================================
\section{Cumprimento dos Objetivos Específicos}
\label{sec:conclusao_objetivos}
% ============================================================

Os objetivos específicos propostos no Capítulo~\ref{cap:introducao} foram cumpridos nas seguintes etapas do desenvolvimento:

\begin{enumerate}
    \item \textbf{Síntese do Controlador PID Linear e Estabilização para Coleta de Dados:} Projetou-se e sintonizou-se um controlador PID clássico na bancada experimental. O controlador desempenhou papel duplo fundamental: garantiu a estabilização da planta em malha fechada durante os ensaios de varredura dinâmica em torno dos equilíbrios operacionais de $45^\circ$ e $60^\circ$, viabilizando a coleta de dados experimentais ricos em amplitude e frequência sem perda de estabilidade, e serviu como linha de base de desempenho linear para a comparação final.
    
    \item \textbf{Identificação Paramétrica do Modelo Dinâmico do Aeropêndulo:} A dinâmica da planta foi identificada por meio do algoritmo FROLS na formulação polinomial NARX, selecionando uma estrutura parcimoniosa de apenas 10 termos a partir de 496 regressores candidatos. O modelo demonstrou estabilidade estrutural e concordância de fase em simulação livre (\textit{free-run}) de longo horizonte. 
    
    \item \textbf{Formulação e Resolução Offline do Problema de Controle Ótimo (NMPC):} O NMPC foi estruturado no ambiente CasADi considerando o modelo preditivo não-linear, restrições operacionais nos atuadores ($0 \le u \le 80\%$) e penalização de esforço e variação de controle ($\Delta u$). A simulação confirmou que o custo computacional de resolução online é incompatível com a taxa de 100\,Hz do sistema mecatrônico, justificando a geração de uma base abrangente de trajetórias ótimas com horizonte de visão antecipada (\textit{preview}) para o aprendizado supervisionado.
    
    \item \textbf{Treinamento da Rede MLP e Implantação em Hardware Embarcado:} A rede neural do tipo MLP ($130 \to 128 \to 128 \to 1$) foi treinada supervisionadamente, alcançando coeficiente de determinação $R^2 = 0{,}999$ e aprendendo a compensar preventivamente a inércia mecânica. A implementação nativa em linguagem C com instruções de ponto flutuante na FPU Cortex-M4 operando a 180\,MHz alcançou tempo de inferência determinístico em torno de 2,07\,ms a 2,08\,ms, consumindo cerca de $20{,}7\%$ da janela de 10\,ms e garantindo uma folga de CPU de $79{,}3\%$. Adicionalmente, a exploração do framework ST X-CUBE-AI com compressão avançada demonstrou a possibilidade de reduzir em $86{,}4\%$ o consumo de memória Flash dos pesos, mantendo o comportamento em malha fechada praticamente indistinguível.
    
    \item \textbf{Comparação Experimental entre Controladores e Fechamento do Gap Sim-to-Real:} Os ensaios espelhados na bancada física frente aos sinais evidenciaram vantagens expressivas da compensação preditiva neural sobre o PID clássico, com redução de erro quadrático médio de rastreamento entre $69\%$ e $81\%$ em manobras agressivas de alta frequência e amplitude. Concomitantemente, a transição para a bancada revelou o viés estacionário e o travamento mecânico por atrito estático de rolamento (\textit{stick-slip}) decorrentes da ausência de ação integradora na MLP pura; a introdução da arquitetura híbrida (MLP com integrador residual paralelo saturado) atenuou o viés em mais de $96\%$ e reduziu o erro senoidal pela metade, fechando eficazmente o \textit{Sim-to-Real gap}.
\end{enumerate}

% ============================================================
\section{Sustentabilidade e Eficiência Energética}
\label{sec:conclusao_sustentabilidade}
% ============================================================

Em consonância com as diretrizes de projeto de engenharia, a solução desenvolvida apresenta impactos positivos diretos sobre eficiência energética e sustentabilidade:

\begin{itemize}
    \item \textbf{Eficiência Computacional e Redução de Consumo Elétrico:} A substituição do solucionador numérico iterativo NMPC (que exige computadores de alto desempenho com consumo típico entre 45\,W e 120\,W) pela inferência de redes neurais em microcontrolador ARM Cortex-M4 permite operar o sistema preditivo com consumo elétrico inferior a 0,5\,W. Essa redução de ordens de magnitude na potência computacional prolonga a autonomia de baterias em veículos não tripulados e viabiliza plataformas móveis energeticamente eficientes.
    \item \textbf{Otimização do Ciclo de Vida de Baterias:} O rastreamento angular mais suave e a atenuação de sobre-sinais (\textit{overshoots}) propiciados pela característica preditiva reduzem picos transitórios de corrente no motor elétrico BLDC, mitigando o estresse eletroquímico e térmico das baterias de polímero de lítio (LiPo) e estendendo sua vida útil operacional.
    \item \textbf{Reaproveitamento de Hardware e Baixo Impacto Material:} O projeto foi concebido sobre uma plataforma de desenvolvimento de baixo custo (STM32 Nucleo-F446RE) e reaproveitou integralmente a estrutura mecatrônica pré-existente da bancada, minimizando o descarte de componentes eletrônicos e o impacto ambiental por geração de lixo eletrônico.
\end{itemize}

% ============================================================
\section{Limitações do Trabalho}
\label{sec:conclusao_limitacoes}
% ============================================================

Embora os resultados experimentais atestem a viabilidade da metodologia desenvolvida, identificam-se aspectos limitantes que delimitam o escopo das conclusões obtidas:

\begin{enumerate}
    \item \textbf{Restrição a Planta com Um Único Grau de Liberdade (1-DOF):} A validação foi restrita ao movimento angular no plano vertical. Em plataformas aéreas completas (como quadrotores e veículos VTOL), manifestam-se acoplamentos inerciais e aerodinâmicos cruzados entre eixos (arfagem, rolamento e guinada), além de momentos giroscópicos decorrentes da variação de rotação de múltiplos motores, efeitos que demandam formulações multivariáveis mais complexas.
    
    \item \textbf{Dependência de Identificação Prévia em Malha Fechada:} A obtenção do modelo dinâmico NARX dependeu da existência prévia de um controlador estabilizador (PID linear) funcional na bancada. Em plantas instáveis em malha aberta para as quais não se dispõe de um controlador preliminar estável, a coleta de dados dinâmica exige metodologias alternativas de síntese ou identificação iterativa.
    
    \item \textbf{Sensibilidade a Discrepâncias Estáticas no Modelo Puramente por Dados:} A rede neural nominal opera como um aproximador de função estático baseado em dados. Discrepâncias entre a planta real e o modelo treinado (tais como perdas aerodinâmicas de recirculação em ângulos elevados e atrito seco de mancais) manifestaram-se como erros persistentes de regime permanente, tornando indispensável a adição de um integrador residual no firmware para assegurar erro estacionário nulo.
    
    \item \textbf{Sensibilidade à Variação de Tensão da Alimentação (Bateria LiPo):} Os parâmetros da rede neural foram sintetizados assumindo uma relação de empuxo constante. Em ensaios operacionais contínuos alimentados por baterias recarregáveis, a curva de descarga reduz a tensão de barramento e diminui o empuxo disponível para uma dada largura de pulso PWM. A rede neural estática não possui mecanismos intrínsecos de compensação para essa perda gradual de ganho do atuador.
    
\end{enumerate}

% ============================================================
\section{Trabalhos Futuros}
\label{sec:conclusao_futuros}
% ============================================================

Como desdobramentos desta pesquisa, sugerem-se as seguintes frentes de investigação:

\begin{enumerate}
    \item \textbf{Extensão para Sistemas com Múltiplos Graus de Liberdade (2-DOF e VTOL):} Aplicar a formulação de modelagem NARX multivariável e controle preditivo neural em bancadas biconduzidas (arfagem e guinada) e veículos aéreos VTOL em voo livre, avaliando o comportamento frente a perturbações atmosféricas e acoplamentos giroscópicos;
    
    \item \textbf{Quantização Inteira Pura em Ponto Fixo (INT8):} Investigar técnicas de quantização pós-treinamento e treinamento ciente de quantização (\textit{Quantization-Aware Training} - QAT) para representação inteira de 8 bits, permitindo a execução da política neural em microcontroladores ultracompactos (como núcleos ARM Cortex-M0+ ou Cortex-M33) desprovidos de unidade de ponto flutuante em hardware;
    
    \item \textbf{Mecanismos de Adaptação Paramétrica Online em Firmware:} Implementar esquemas adaptativos em tempo real, tais como estimadores recursivos de ganho de atuador baseados em Filtro de Kalman ou aprendizado por reforço residual, capazes de compensar a degradação de tensão da bateria e variações de carga útil ao longo do tempo de operação;
    
    \item \textbf{Certificação Formal de Estabilidade e Limites de Lipschitz:} Desenvolver métodos de análise baseados na constante de Lipschitz da rede treinada ou em funções de controle de barreira (\textit{Control Barrier Functions} - CBF) para delimitar o envelope seguro de operação do controlador neural embarcado.
\end{enumerate}